\VerdierPhysicalPageAnchor{0050}%
\VerdierSemanticAnchor{ch1-1-4-7-continuation}
where the $(F_i,p_i)$ are the twisted $G$-functors determined
respectively by the functors $(F_i,m_i)$ and the sections $s_i$, are
$\beta_{j,i}$-morphisms, where

\VerdierSemanticAnchor{ch1-1-4-7-eq3}
\begin{equation}
  \beta_{j,i}(\sigma)=s_j(\sigma)s_i(\sigma)^{-1},
  \qquad i,j\in\{1,2,3\},\quad i<j.
  \tag{1.4.7.3}
\end{equation}

Thus we have the relation

\VerdierSemanticAnchor{ch1-1-4-7-eq4}
\begin{equation}
  \beta_{3,1}=\beta_{3,2}\beta_{2,1}.
  \tag{1.4.7.4}
\end{equation}

\VerdierSemanticAnchor{ch1-1-4-7-composition}
Relations (1.4.7.2) and (1.4.7.4) then make it possible to define the
\emph{composite} of the $\beta_{3,2}$-morphism $v$ with the
$\beta_{2,1}$-morphism $u$, and the composition so defined is
compatible with the composition of morphisms of $H$-functors.

\VerdierSemanticAnchor{ch1-1-4-8}
\noindent\textbf{1.4.8.}\quad
Let $\mathcal C'$ be an additive category and let
\[
  (F,p)\colon(\mathcal C,G,T,c)\longrightarrow
  (\mathcal C',G,T',c')
\]
be a $G$-functor twisted by a normalized cocycle
$\alpha\in Z^2(G,\mathbb Z/2\mathbb Z)$. We denote by

\VerdierSemanticAnchor{ch1-1-4-8-eq1}
\begin{equation}
  (F^{\circ},p^{\circ})\colon
  (\mathcal C^{\circ},G^{\circ},T^{\circ},c^{\circ})
  \longrightarrow
  (\mathcal C'^{\circ},G^{\circ},T'^{\circ},c'^{\circ})
  \tag{1.4.8.1}
\end{equation}
the \emph{opposite twisted graded functor} to the functor $(F,p)$.
(We use the notation of (1.3.9).) It is defined as follows.

\VerdierSemanticAnchor{ch1-1-4-8-opposite-data}
The underlying functor is $F^{\circ}$, the functor opposite to $F$.
For every $g\in G^{\circ}$ (the group opposite to $G$), the
isomorphism in $\mathcal C^{\circ}$
\[
  p^{\circ}(g)\colon
  F^{\circ}T^{\circ}(g)=FT(g^{-1})
  \overset{\sim}{\longrightarrow}
  T'^{\circ}(g)F^{\circ}=T'(g^{-1})F
\]
is equal to $p^{-1}(g^{-1})$.

\VerdierSemanticAnchor{ch1-1-4-8-opposite-cocycle}
The functor $(F^{\circ},p^{\circ})$ is a $G^{\circ}$-functor twisted
by the cocycle $\alpha^{\circ}$ in
$Z^2(G^{\circ},\mathbb Z/2\mathbb Z)$ opposite to the cocycle
$\alpha$:
\[
  \alpha^{\circ}(g,g')=\alpha^{-1}(g'^{-1},g^{-1}).
\]

\VerdierSemanticAnchor{ch1-1-4-9}
\noindent\textbf{1.4.9.}\quad
The only property of the group $\mathbb Z/2\mathbb Z$ used in this
paragraph is that it acts canonically on the identity functor of
additive categories. One can therefore develop analogous
considerations by replacing the group $\mathbb Z/2\mathbb Z$ with a
commutative group $K$ acting on the identity functors of the categories
under consideration, compatibly with the operations of $G$.
